\section{Formation energies}\label{sec:formation}

Clusters are put in context by their formation energies: how much a
structure's energy differs from a stoichiometric mix of isolated,
single-component reference structures. A cluster database already contains
its own references, namely the monomers. \CJKQC{} scans the run's final
records for monomers and prints each printed record's formation energy and
formation free energy as a section after the main table.

Nothing is stored. Formation values are derived at request time from the
records' stored values under the run's thermochemistry conditions, so
temperature, reference frequency, and unit changes apply to them like to any
other derived property.

The sign convention is:

\begin{quote}
\textbf{negative} = more stable than the isolated monomers.
\end{quote}

The binding energy is the same number with the opposite sign; there is no
separate binding column.

\subsection{Monomer references}\label{sec:monomer-references}

A \textbf{monomer record} is an \fld{ok} record whose composition is exactly
one component with count 1. Thus \texttt{\{"am": 1\}} qualifies, while
\texttt{\{"am": 2\}} and \texttt{\{"am": 1, "sa": 1\}} do not. Records
without components never qualify.

Monomers are collected from exactly the records the run prints: the records
that remain after every filtering step in a filtering run, or every loaded
record in a print-only run. A monomer that a filter removes is gone:
it prints no row and serves as no reference.

For each component code, the reference is a single monomer record: the first
of that code's monomers under the run's sort ordering.

\begin{center}
\begin{tabular}{@{}ll@{}}
\toprule
\textbf{Sort} & \textbf{Reference monomer} \\
\midrule
\fld{el}, or no sort & lowest electronic energy \\
\fld{g} & lowest Gibbs free energy at the run's conditions \\
\fld{rg} & first under radius of gyration \\
\fld{mass} & first under total mass \\
\fld{filename} & first under the file-stem ordering (nameless first) \\
\bottomrule
\end{tabular}
\end{center}

The ordering is always ascending; \flg{--reverse} orients the table, never
the reference. Ties keep the first in load order. For numeric properties, a
monomer whose sorting value is undefined orders after every defined one.

Both reference values come from that one record:

\begin{itemize}
  \item the \textbf{electronic reference} is its \fld{electronic\_energy};
  \item the \textbf{Gibbs reference} is its Gibbs free energy derived at the
        run's conditions.
\end{itemize}

Under \flg{--sort g}, the lowest-G conformer supplies the electronic
reference too, even when another conformer's electronic energy is lower. A
level whose value the chosen record lacks has no reference at that level. A
code with no qualifying monomer has no reference.

\subsection{The raw formation section}\label{sec:raw-formation}

\flg{--formation} is a print flag: it produces output and satisfies the
visible-outcome rule on its own, but adds no column to the normal table. The
section prints after the main table (if one was requested) and before any
summary line.

For a printed record with components $C = \{\text{code}_i : n_i\}$:

\begin{align}
E_{\mathrm{form}} &= E_{\mathrm{el}}
  - \sum_i n_i\,E_{\mathrm{ref}}(\text{code}_i),\\
G_{\mathrm{form}} &= G - \sum_i n_i\,G_{\mathrm{ref}}(\text{code}_i).
\end{align}

Both reference values are the ones selected above. A level is undefined when
the record lacks \fld{components}, lacks the level's input energy (electronic
energy for the electronic level, defined thermochemistry for the Gibbs
level), or any referenced code lacks that level's monomer reference.
Undefined is a property of the level: a record can have a defined formation
energy and an undefined formation free energy, and vice versa.

The section is a table with header \fld{name el g}, one row per printed
record in run order. \fld{name} is the short composition label (for example
\texttt{2am3sa}), or \texttt{-} when the record has no components.
\fld{el} and \fld{g} are the formation energy and formation free energy in
the run's energy unit: hartree by default, kcal/mol under \flg{--unit}. An
undefined value prints \texttt{-}.

\subsection{The corrected formation section}\label{sec:corrected-formation}

\flg{--formation-corr} prints a second section evaluated at the attached
single-point corrections. Every input energy is a corrected value; a record
with no correction has no values at all, so corrected and raw levels are
never mixed.

The reference monomers are the same set as for the raw section, selected the
same way. A code's corrected electronic reference is that monomer's
correction energy, and its corrected Gibbs reference is the monomer's
\fld{gcorr} value. A level's reference is undefined when the selected
monomer carries no correction, or its corrected Gibbs is undefined.

For a printed record:

\begin{align}
E_{\mathrm{form,corr}} &= E_{\mathrm{corr}}
  - \sum_i n_i\,E_{\mathrm{corr,ref}}(\text{code}_i),\\
G_{\mathrm{form,corr}} &= G_{\mathrm{corr}}
  - \sum_i n_i\,G_{\mathrm{corr,ref}}(\text{code}_i).
\end{align}

The section has header \fld{name elcorr gcorr}. It prints after the raw
formation section, when one was requested. A run with
\flg{--formation-corr} but no other print flag prints only the corrected
section.

\subsection{Concentrations and \texttt{-{}-cnt}}\label{sec:concentrations}

\flg{--conc CODE VALUE} evaluates the formation free energies at a vapor
concentration different from the standard state. It may be repeated, with
one concentration setting per code. The code must be a non-empty, letter-only
component code. When a formation section is printed, every such code must
appear as a component of at least one record in the section's record set.
Specifying the same code twice is a usage error.

VALUE is a partial pressure in pascals, optionally suffixed:

\begin{center}
\begin{tabular}{@{}lll@{}}
\toprule
\textbf{Form} & \textbf{Meaning} & \textbf{Pressure in Pa} \\
\midrule
\texttt{100} & pascals & $100$ \\
\texttt{100ppt} & parts per trillion by mole fraction & $100\cdot10^{-12}p^\circ$ \\
\texttt{10ppb} & parts per billion by mole fraction & $10\cdot10^{-9}p^\circ$ \\
\texttt{1e10cmmc} & molecules per cm$^3$ & $c\cdot10^{6}k_B T$ \\
\bottomrule
\end{tabular}
\end{center}

Here $p^\circ = 101325$ Pa. The value must be finite and positive; the
\texttt{\^{}} power syntax is not part of the grammar. For powers of ten, use
scientific notation instead, for example \texttt{1e10} or \texttt{1E10} for
$10^{10}$. The concentration term
is a chemical-potential (entropy) correction, so it shifts only the Gibbs
levels: the raw section's \fld{g} and the corrected section's \fld{gcorr}.
The electronic levels \fld{el} and \fld{elcorr} stay at the standard state.

For the set of concentration settings $S$ and a record with components $C$:

\begin{equation}
g_{\mathrm{conc}} = g_{\mathrm{form}}
  - \sum_{\text{code}_i \in S} n_{i,\mathrm{eff}}\,
    \frac{R T \ln(c_i/p^\circ)}{N_A E_h},
\end{equation}

where $c_i$ is the resolved partial pressure and $n_{i,\mathrm{eff}} = n_i$.
Under \flg{--cnt}, the self-consistent form is used:

\begin{equation}
n_{i,\mathrm{eff}} = n_i\,\frac{N-1}{N},
\qquad N = \sum_i n_i .
\end{equation}

The self-consistent form keeps the concentration term of a monomer at zero
under any concentration, so a monomer's formation free energy stays exactly
zero while a non-reference conformer keeps its conformer difference. A
concentration setting whose code is absent from the record's components
contributes nothing.

\subsection{Boltzmann-averaged sections}\label{sec:bavg}

\flg{--bavg} collapses each printed formation section into one
Boltzmann-averaged row per composition. It is a settings flag: it configures
the section's presentation, never counts as a manipulation or print flag, and
is an accepted no-op without a printed formation section. The main table is
never affected; it still prints one row per record.

The collapse groups the section's record set by \fld{components}, always by
composition (\flg{--noname} never pools the section), in first-appearance
order. Records without components form one unnamed ensemble. A group's
members are its \fld{ok} records that define the section's Gibbs level: the
derived Gibbs for the raw section, \fld{gcorr} for the corrected one. Failed
and truncated records contribute nothing, and an \fld{ok} record without the
level value contributes nothing.

An ensemble with members carries two values:

\begin{itemize}
  \item the \textbf{electronic value} is the members' population-weighted
        mean $\sum_j p_j e_j$, with the members' Boltzmann populations at the
        run's temperature. Here $e_j$ is \fld{electronic\_energy} for the raw
        section and the correction energy for the corrected section;
  \item the \textbf{Gibbs value} is the ensemble free energy
        \begin{equation}
        G_{\mathrm{ens}} = -k_B T \ln\!\left(\sum_j e^{-g_j/(k_B T)}\right)
        = g_{\min} - k_B T \ln Z,
        \end{equation}
        computed against the members' minimum as in the population
        numerics. This keeps the conformational entropy term
        $-k_B T\ln Z$; a plain weighted mean would drop it.
\end{itemize}

The two values are defined together or not at all: a member that defines the
Gibbs level always defines the electronic level, because the Gibbs derives
from the electronic energy.

Over the collapsed set, the standing rules apply unchanged. An ensemble whose
composition is exactly one component with count 1 is its code's monomer
ensemble and supplies both reference values; a monomer row's formation values
are therefore exactly zero. An ensemble without values supplies no reference.

The section prints one row per ensemble in first-appearance order, using the
same header, unit, and formatting as the uncollapsed section. A composition
whose records all lack the section's Gibbs level prints one row of dashes.
Undefined members never poison their group: the surviving members define the
ensemble.

\subsection{A formation example}\label{sec:formation-example}

Print formation energies and free energies for the lowest conformer of each
composition, using the lowest-Gibbs monomer of each code as the reference:

\begin{verbatim}
CJKQC db.jsonl --sort g --glob --formation --unit \
  --filename --name --g
\end{verbatim}

Evaluate the same formation free energies at an ammonia vapor of
$10^{-9}$ bar-equivalent mole fraction, with the self-consistent form:

\begin{verbatim}
CJKQC db.jsonl --formation --conc am 10ppb --cnt --unit --name
\end{verbatim}

Collapse the section into one Boltzmann-averaged row per composition:

\begin{verbatim}
CJKQC db.jsonl --formation --bavg --unit --name
\end{verbatim}

In the last example the main table, if requested, would still show one row
per record; only the formation section is collapsed.
